== 3.4 MATRIX OPERATIONS ==

Matrix Addition (same size only, add entry by entry):
$$A + B = [a_{ij} + b_{ij}]$$

Example: [1 2 / 3 4] + [5 6 / 7 8] = [6 8 / 10 12]

Properties:
$$A+B = B+A$$
$$(A+B)+C = A+(B+C)$$

Scalar Multiplication (multiply every entry by c):
$$cA = [c \cdot a_{ij}]$$

Matrix Multiplication (A is m x n, B is n x p):
$$(AB)_{ij} = \sum_{k=1}^{n} a_{ik} b_{kj}$$

Example: A=[1 2/3 4], B=[0 1/2 3]
Row1*Col1: (1)(0)+(2)(2)=4
Row1*Col2: (1)(1)+(2)(3)=7
Row2*Col1: (3)(0)+(4)(2)=8
Row2*Col2: (3)(1)+(4)(3)=15
Result: AB=[4 7/8 15]

WARNING: AB is generally NOT equal to BA

Properties:
$$(AB)C = A(BC)$$
$$A(B+C) = AB+AC$$
$$AI = IA = A$$


== 3.5 MATRIX INVERSES ==

2x2 Inverse Formula:
If A=[a b/c d] and det(A)=ad-bc is not 0:
$$A^{-1} = \frac{1}{ad-bc}\begin{pmatrix}d & -b \\ -c & a\end{pmatrix}$$

Example: A=[2 1/5 3], det=(2)(3)-(1)(5)=1
$$A^{-1} = [3\ -1\ /\ -5\ 2]$$

Check: A times A^{-1} = identity matrix.

Row Reduction Method (works for any size):
Form [A|I] and row reduce until left side = I.
Right side becomes A^{-1}.

Example: A=[1 2 0/2 5 0/3 7 1]
R2=R2-2R1: [0 1 0|-2 1 0]
R3=R3-3R1: [0 1 1|-3 0 1]
R3=R3-R2:  [0 0 1|-1 -1 1]
R1=R1-2R2: [1 0 0|5 -2 0]
Result: A^{-1}=[5 -2 0/-2 1 0/-1 -1 1]

Theorem 7 (all equivalent for n x n matrix A):
$$A \text{ invertible} \Leftrightarrow Ax=0 \text{ has only trivial soln}$$
$$\Leftrightarrow \text{RREF of A is } I_n$$
$$\Leftrightarrow A \text{ has n pivot positions}$$
$$\Leftrightarrow Ax=b \text{ has exactly one solution}$$
$$\Leftrightarrow \text{columns linearly independent}$$
$$\Leftrightarrow \det(A) \neq 0$$


== 3.6 DETERMINANTS ==

2x2 Formula:
$$\det\begin{pmatrix}a & b \\ c & d\end{pmatrix} = ad - bc$$

Example:
$$\det\begin{pmatrix}3 & 2 \\ 1 & 4\end{pmatrix} = (3)(4)-(2)(1) = 10$$

Cofactor Expansion:
$$C_{ij} = (-1)^{i+j} M_{ij}$$

M_ij = det with row i and col j deleted.

Sign pattern (+/-) for 3x3:
$$\begin{pmatrix}+ & - & + \\ - & + & - \\ + & - & +\end{pmatrix}$$

Expanding along row i:
$$\det(A) = a_{i1}C_{i1} + a_{i2}C_{i2} + a_{i3}C_{i3}$$

TIP: pick the row or col with the most zeros!

Example (expand along row 1):
A = [1 2 3 / 0 4 5 / 0 0 6]
$$= 1\cdot\det\begin{pmatrix}4&5\\0&6\end{pmatrix} - 2\cdot\det\begin{pmatrix}0&5\\0&6\end{pmatrix} + 3\cdot\det\begin{pmatrix}0&4\\0&0\end{pmatrix}$$
$$= 1(24) - 2(0) + 3(0) = 24$$

Shortcut: triangular matrix det = product of diagonals.

Row Operation Effects on det:
$$R_i \leftrightarrow R_j \Rightarrow \text{det changes sign}$$
$$R_i \to cR_i \Rightarrow \text{det multiplied by } c$$
$$R_i \to R_i + cR_j \Rightarrow \text{det unchanged}$$

Theorem 2:
$$A \text{ invertible} \Leftrightarrow \det(A) \neq 0$$
$$\det(AB) = \det(A)\det(B)$$
$$\det(A^{-1}) = \frac{1}{\det(A)}$$
$$\det(A^T) = \det(A)$$


== 4.2 VECTOR SPACES ==

A vector space V satisfies 10 axioms (closure, identity,
inverse, distributive, associative laws, etc.)
Examples: R^n, all m x n matrices, polynomials, functions.

Subspace Test (H is a subspace of V if ALL 3 hold):
1. Zero vector 0 is in H
2. u,v in H implies u+v in H (closed under addition)
3. u in H, c in R implies cu in H (closed under scalar mult)

Example (IS a subspace): H={(x,y): y=3x} in R^2
$$0: 0=3(0)\ \checkmark$$
$$(x_1+x_2, 3x_1+3x_2) = (x_1+x_2, 3(x_1+x_2))\ \checkmark$$
$$c(x,3x)=(cx,3cx)\ \checkmark$$

Example (NOT a subspace): H={(x,y): y=x+1}
(0,0) not in H. Fails test 1. Not a subspace.

Solution Subspaces:
The solution set of Ax=0 is always a subspace (= Nul(A)).

Example: A=[1 -2/2 -4]
Row reduce: [1 -2/0 0], so x1=2x2, x2 free.
$$\text{Solution: } t\begin{pmatrix}2\\1\end{pmatrix}, t \in \mathbb{R}$$
This is a line through origin, which is a subspace.


== 4.3 SPAN AND LINEAR INDEPENDENCE ==

Span:
$$\text{Span}\{v_1,...,v_p\} = \{c_1v_1+...+c_pv_p : c_i \in \mathbb{R}\}$$

This is always a subspace.

Linear Independence:
{v1,...,vp} is linearly independent if:
$$c_1v_1+...+c_pv_p = 0 \Rightarrow \text{all } c_i = 0$$

Test: form matrix with vi as columns and row reduce.
Every column has a pivot => independent.
Any free column => dependent.

Example (dependent): v1=[1/2/3], v2=[2/4/6]
v2 = 2*v1, so dependent (2v1 - v2 = 0).

Example (independent): v1=[1/0/0], v2=[0/1/0], v3=[0/0/1]
Matrix = I_3. All pivots. Independent.

Key Facts:
More vectors than entries in each => always dependent.
Set containing 0 => always dependent.
Single nonzero vector => always independent.


== 4.4 BASIS AND DIMENSION ==

A basis for subspace H is a set {v1,...,vp} that is:
(1) linearly independent  AND  (2) spans H.

Standard basis for R^n: {e1, e2, ..., en}

Finding basis for Col(A):
Row reduce A. Pivot columns of ORIGINAL A form a basis.

Example: A=[1 2 0/0 0 1/0 0 0]
Pivots in columns 1 and 3.
$$\text{Basis for Col}(A) = \left\{\begin{pmatrix}1\\0\\0\end{pmatrix}, \begin{pmatrix}0\\1\\0\end{pmatrix}\right\}$$

Example 2: Is {[1/0/0],[1/1/0],[1/1/1]} a basis for R^3?
$$\det = 1 \neq 0 \Rightarrow \text{Yes, it is a basis.}$$

Dimension:
$$\dim(H) = \text{number of vectors in any basis for } H$$
$$\dim(R^n) = n, \quad \dim(\{0\}) = 0$$

Basis for Solution Space:
Row reduce A. Express basic vars in terms of free vars.
Each free variable gives one basis vector.

Example: A=[1 2 -1 0/0 0 0 0]
Free: x2, x3. Basic: x1 = -2x2 + x3.
$$x = x_2\begin{pmatrix}-2\\1\\0\end{pmatrix} + x_3\begin{pmatrix}1\\0\\1\end{pmatrix}$$
$$\text{Basis: } \{[-2,1,0]^T,\ [1,0,1]^T\}, \quad \dim = 2$$


== 4.5 MATRIX SPACES ==

Null Space of A (m x n matrix):
$$\text{Nul}(A) = \{x : Ax = 0\}$$

Subspace of R^n. Find by row reducing [A|0].

Column Space of A:
$$\text{Col}(A) = \text{Span of columns of } A$$

Subspace of R^m.
Basis = pivot columns of ORIGINAL A (not RREF!).

Rank:
$$\text{rank}(A) = \dim(\text{Col}(A)) = \text{number of pivot columns}$$

Rank-Nullity Theorem (A is m x n):
$$\text{rank}(A) + \text{nullity}(A) = n$$

nullity(A) = dim(Nul(A)) = number of free variables.

Example 1: A is 3x5, rank(A)=2
$$\text{nullity}(A) = 5 - 2 = 3$$

Example 2: A=[1 0 -2/0 1 3]
n=3, rank=2, nullity=1.
Free: x3. Solve: x1=2x3, x2=-3x3.
$$\text{Basis for Nul}(A): \{[2,-3,1]^T\}$$
$$\text{Basis for Col}(A): \{[1,0]^T, [0,1]^T\}$$


== 5.1 SECOND-ORDER LINEAR DEs ==

General Form:
$$ay'' + by' + cy = f(t)$$

Homogeneous when f(t)=0:
$$ay'' + by' + cy = 0$$

Theorem 1 (Superposition):
If y1 and y2 solve the homogeneous equation, then:
$$y = c_1y_1 + c_2y_2 \text{ is also a solution}$$

Theorem 4 (General Solution):
If y1, y2 are linearly independent solutions:
$$y = c_1y_1 + c_2y_2 \text{ is the GENERAL solution}$$

Wronskian (check linear independence):
$$W(y_1,y_2) = y_1y_2' - y_2y_1' \neq 0$$

Characteristic Equation Method:
Substitute y=e^{rt} into ay''+by'+cy=0 to get:
$$ar^2 + br + c = 0$$

--- CASE 1: Two distinct real roots r1 != r2 ---
$$y = c_1e^{r_1 t} + c_2e^{r_2 t}$$

Example: y'' - 5y' + 6y = 0
$$r^2 - 5r + 6 = (r-2)(r-3) = 0$$
$$r_1=2,\ r_2=3$$
$$y = c_1e^{2t} + c_2e^{3t}$$

--- CASE 2: Repeated root r1=r2=r ---
$$y = (c_1 + c_2t)e^{rt}$$

Example: y'' - 4y' + 4y = 0
$$r^2 - 4r + 4 = (r-2)^2 = 0 \Rightarrow r=2$$
$$y = (c_1 + c_2t)e^{2t}$$

--- CASE 3: Complex roots r = alpha +/- beta*i ---
$$y = e^{\alpha t}(c_1\cos(\beta t) + c_2\sin(\beta t))$$

Example 1: y'' + 4y = 0
$$r^2 + 4 = 0 \Rightarrow r = \pm 2i$$
$$\alpha=0,\ \beta=2$$
$$y = c_1\cos(2t) + c_2\sin(2t)$$

Example 2: y'' - 2y' + 5y = 0
$$r = \frac{2 \pm \sqrt{4-20}}{2} = \frac{2 \pm 4i}{2} = 1 \pm 2i$$
$$\alpha=1,\ \beta=2$$
$$y = e^{t}(c_1\cos(2t) + c_2\sin(2t))$$

--- Initial Value Problems ---
Write general y and y'. Plug in t0. Solve 2x2 system for c1, c2.

Example: y'' - y = 0, y(0)=2, y'(0)=1
$$r^2 - 1 = 0 \Rightarrow r = \pm 1$$
$$y = c_1e^{t} + c_2e^{-t}$$
$$y' = c_1e^{t} - c_2e^{-t}$$
At t=0:
$$c_1 + c_2 = 2 \quad \text{and} \quad c_1 - c_2 = 1$$
$$c_1 = \frac{3}{2},\ c_2 = \frac{1}{2}$$
$$y = \frac{3}{2}e^{t} + \frac{1}{2}e^{-t}$$
